\documentclass[10pt]{article}
\usepackage[margin=20mm]{geometry}
\usepackage{amsmath,amssymb,fontspec,unicode-math,xurl,hyperref}
\IfFontExistsTF{DejaVu Serif}{\setmainfont{DejaVu Serif}}{\setmainfont{Latin Modern Roman}}
\setmathfont{latinmodern-math.otf}
\setlength{\parindent}{0pt}\setlength{\parskip}{6pt}\sloppy
\title{Technical Verification of
Gradient-Free Cutting-Plane Active Learning}
\author{AI-assisted technical review}
\date{Initial review, 25 October 2026}
\begin{document}
\maketitle
Abstract: The assessment is Major. The proof of prediction convergence incorrectly converts volume decay into point identification. The inexact-cut bound is false on an explicit full-dimensional ReLU parameter domain, and its dense-data corollary fails even from the specified unit-ball initialization. Fixed unit margins conflict with the claimed scale-invariant radius, while the printed strict-mistake rule cannot leave the zero center. The query counter counts accepted cuts instead of all acquired labels. Two released-code decoding defects are Major for the frozen three-layer algorithm and sentiment evaluation paths; their historical experimental impact remains unresolved. The valid conditional centroid-cut volume bound and the inspected two-layer convex reformulation are distinguished from these failures.\par
\section*{Identity, version, and review scope}
The paper is Active Learning of Deep Neural Networks via Gradient-Free Cutting Planes, by Erica Zhang, Fangzhao Zhang, and Mert Pilanci, ICML 2025, PMLR 267; I25-43666, display 0297, OpenReview `vZhdXRnfPU` [paper]. The reviewed source is the staged 49-page proceedings PDF with embedded Appendices A–H. Page numbers below are PDF pages and agree with printed numbering.\par
The review checked the two-layer equivalence, the three-layer construction, the definitions and algorithms underlying Sections 5–6, the centroid and inexact-cut convergence arguments, experimental design and reported scalar comparisons, and released core/source application paths. The general-depth indexing and all external approximation results were not independently certified. Exact raw historical query logs and results were unavailable.\par
The released repository `pilancilab/cpal` was inspected at commit `14c27c1a97756493738390775c1e4fe818d52935` [code]. Its README identifies it as preliminary. Forty selected source members were checked against archive bytes and Git blob identities. No author program was executed. This version uses analytic centers rather than the center of gravity used in the principal theorem, treats a zero prediction differently from the printed pseudocode, and uses zero-margin classification cuts. Those differences are material counterevidence against equating written-algorithm failures with historical experimental failures.\par
\section*{Necessary definitions and valid components}
With fixed activation patterns \(D_i\) and variables \(\theta=(u'_1,u_1,\ldots,u'_P,u_P)\), the lifted two-layer score is linear: \[f_\theta(X)=\sum_iD_iX(u'_i-u_i).\] The signs of the primed and unprimed variables are interchanged between some earlier and later displays; the examples use the Section 6 convention and are invariant to consistently swapping these roles. Activation constraints require \((2D_i-I)Xu'_i\ge0\) and \((2D_i-I)Xu_i\ge0\).\par
Theorem 4.1’s splitting of positive and negative output weights, merging inside each activation cone, and reconstruction with at most \(2P\) neurons were checked. The three-layer construction similarly duplicates signed branches under the stated ample-width condition. These are feasibility reformulations; they do not by themselves supply an initial radius, a statistically meaningful label budget, or prediction identification.\par
For a full-dimensional convex body, an active valid linear cut excluding its centroid retains at most \(1-1/e\) of the volume. Adding activation constraints cannot increase that retained volume. This portion of Theorem 6.3 is sound within its assumptions. It is a per-active-cut statement, not a guarantee for arbitrary queries, arbitrary analytic centers, or convergence to one prescribed real-valued prediction.\par
\section*{Problems in the central guarantees and algorithms}
Problem 1 (Volume decay is incorrectly promoted to prediction identification).  Assessment: Major.\par
Appendix G.1, p. 30, asserts that nestedness, nonemptiness, and vanishing volume imply \(\bigcap_tT^{(t)}=\{\theta^*\}\). Appendix G.2, pp. 30–31, then uses the purported implication \(\theta^{(k)}\to\theta^*\) to prove Corollary 6.4’s convergence of real-valued scores on every fixed \(X\).\par
The geometric implication is false even for active homogeneous cuts. Let \(T_k\) be the closed unit-disk sector with angles \(0\le\phi\le\alpha_k\), where \(\alpha_k=(\pi/2)3^{-k}\). Each is a compact convex body with nonempty interior. Cutting at angle \(\alpha_k/3\) retains \(T_{k+1}\), removes the old centroid (whose angle is \(\alpha_k/2\)), and reduces area by exactly \(1/3\). Nevertheless, \[\bigcap_kT_k=[0,1]\times\{0\},\qquad
 \operatorname{CG}(T_k)=
 \left(\frac{2\sin\alpha_k}{3\alpha_k},
       \frac{2(1-\cos\alpha_k)}{3\alpha_k}\right)
 \longrightarrow(2/3,0).\] The retained feasible point \((1/3,0)\) is not the limit, and the linear output equal to the first coordinate differs by \(1/3\) in the limit. This is an exact counterexample to the geometric inference used by the proof; it is not asserted to be a complete network-query trajectory.\par
The network setting also preserves scale ambiguity: zero-margin classification and activation constraints are homogeneous, so a feasible nonzero ray segment can remain inside the initial ball. Correct classification does not fix score amplitude. The finite-intersection property guarantees a nonempty intersection, not a singleton. Corollary 6.4 is therefore unproved by the supplied argument, directly affecting a central prediction-convergence claim. No failure of the conditional volume inequality itself is alleged.\par
A precise sufficient additional assumption would be \(\operatorname{diam}(T_k)\to0\) with \(\theta^*\in T_k\). For a fixed lifted matrix \(A_X\), it then gives \(\|A_X\theta^{(k)}-A_X\theta^*\|\le
\|A_X\|\operatorname{diam}(T_k)\to0\). This conditional repair is valid; the proposed algorithm has not been shown to enforce its new premise.\par
Problem 2 (The inexact-cut bounds fail on admissible ReLU domains).  Assessment: Major.\par
\subsection*{Theorem G.2: explicit counterexample for every proposed dimension constant}
Theorem G.2, pp. 32–34, claims \[\frac{\operatorname{vol}(T_1)}{\operatorname{vol}(T)}
 <1-e^{-1}(1-\widetilde\lambda)^d,\qquad
 \widetilde\lambda=
 -\frac{R^dH_{T^-}}{\Theta_d\operatorname{vol}(T)H_{T^+}}h,\] where \(h\) is the centroid’s distance to the cut, \(R\) is the radius of a \((d-1)\)-ball with the cross-section’s volume, and \(\Theta_d>0\) depends only on dimension.\par
Use input dimension one, \(X=[1]\), label \(+1\), and the complete pattern set \(\{1,0\}\). The lifted parameter dimension is four. In orthogonal coordinates set \[q=(u'_1-u_1)/\sqrt2,\qquad s=(u'_1+u_1)/\sqrt2,\] and define the rectangular body \[-\epsilon\le q\le1,\quad 2\le s\le3,\quad
 -2\le u'_0,u_0\le-1,\qquad 0<\epsilon\le0.1.\] It is compact and full-dimensional. Both active-pattern variables are positive because \(s>|q|\), and both inactive variables are negative. Thus every activation constraint is satisfied strictly throughout \(T\). The classification cut is exactly \(q\ge0\) and removes no additional volume through activation constraints.\par
Its volume is \(1+\epsilon\), its cross-section at \(q=0\) has three-volume one, and \[R=(3/(4\pi))^{1/3},\quad
 h=(1-\epsilon)/2,\quad
 H_{T^-}/H_{T^+}=-\epsilon.\] Writing \(K=R^4/(2\Theta_4)>0\), the printed expression gives \(\widetilde\lambda=K\epsilon(1-\epsilon)/(1+\epsilon)\). For any positive \(\Theta_4\), choose \(\epsilon=\min\{0.1,1/(10K)\}\). Then \[\frac{\operatorname{vol}(T_1)}{\operatorname{vol}(T)}
   =\frac1{1+\epsilon}\ge\frac{10}{11},
 \qquad
 1-e^{-1}(1-\widetilde\lambda)^4
   \le1-e^{-1}(0.9)^4=0.7586343\ldots .\] The claimed upper bound is contradicted. All quantities are well-defined, both half-spaces meet \(T\), and \(0<\widetilde\lambda\le0.1\). Uniformly scaling \(T\) by \(1/5\) places it inside the unit ball, preserves all signs and the retained fraction, and multiplies \(\widetilde\lambda\) by \(1/5\), strengthening the contradiction. This refutes the theorem’s stated universal convex-body bound; the box is not claimed to be a particular recorded historical iterate.\par
\subsection*{Corollary G.3: failure from the algorithm’s actual initialization}
Corollary G.3, p. 35, claims that uniform, increasingly dense input data and minimum raw margin recover the \((1-1/e)^t\) volume bound for Algorithm 5. Take independent inputs uniform on \((0,1)\) and give every input label \(+1\). Both labels being represented is not assumed. The full activation-pattern set remains \(\{1,0\}\); a strictly correct four-neuron representation exists.\par
Starting at \(B_2\subset\mathbb R^4\), every positive input induces the identical constraints \[u'_1\ge u_1\ge0,\qquad u'_0\le0,\quad u_0\le0.\] The first cut retains one orthant, with half of its active-coordinate ordering, so \(\operatorname{vol}(T^{(1)})/\operatorname{vol}(T^{(0)})=1/32\). All later distinct queries add redundant copies of the same constraints. Hence the ratio remains \(1/32\) for every \(t\ge1\), independently of dataset size or minimum-margin selection. Already at \(t=8\), \[(1-1/e)^8=0.02549173\ldots < 1/32=0.03125.\] This is a direct counterexample to the claimed asymptotic bound using the specified initialization and valid ReLU cuts.\par
The proof also confuses raw margin with normalized distance: \(h=|f_\theta(x)|/\|a(x)\|\). Here \(a(x)=x(1,-1,0,0)\). Small positive \(x\) makes the raw margin small but leaves the hyperplane and its distance unchanged. Dense inputs do not guarantee informative cuts. Theorem 6.3 survives because these later cuts are not active. No complete corrected inexact-cut theorem is certified.\par
Problem 3 (The fixed unit margin is incompatible with a universal unit-ball radius).  Assessment: Major.\par
Section 5, p. 5, claims the unit ball contains optimal parameters by scale invariance, while Algorithm 1’s cut enforces margin at least one, as in Equation (3). Scaling preserves the sign, but not that fixed margin.\par
The incompatibility persists with explicit bias augmentation and both classes. Let \(x_+=(0.1,1)\), \(x_-=(-0.1,1)\) and labels \(+1,-1\). These independent rows realize all four first-layer patterns, so \(P=4\) and the stated ample-width condition is met by \(m=8\). The original network \[f(z)=10(z_1)_+-10(-z_1)_+\] (padded with zero neurons) has both margins exactly one.\par
For any lifted parameter with \(\|\theta\|_2\le1\), the reconstructed ReLU function is Lipschitz with constant at most \[\sum_{i=1}^{P}(\|u_i\|_2+\|u'_i\|_2)
 \le\sqrt{2P}\|\theta\|_2\le\sqrt8.\] The two outputs can therefore differ by at most \(0.2\sqrt8=0.565685\ldots\), whereas the two unit margins require a difference of at least two. Thus the radius excludes every solution of a feasible original instance.\par
This invalidates the initialization argument for the central fixed-margin training scheme. A data-dependent adequate radius, or an explicitly changed margin/domain, is required. The zero-margin cuts in Algorithms 2/4 and in the release avoid this particular contradiction; that distinction is not a repair of Algorithm 1 as printed.\par
Problem 4 (The printed strict-mistake rule stalls at the initial center).  Assessment: Major; written algorithm, with different released behavior.\par
Algorithms 1, 2, and 4 (pp. 4, 6, 16) start from the unit ball and cut only if \(y f_{\theta_c}(x)<0\). Its centroid and symmetric barrier center are zero, and \(f_0(x)=0\) for every input. Therefore no label, including either binary class, can trigger the strict inequality. The localization set and center remain unchanged by induction. In Algorithms 2/4, \(D_{\rm AL}\) stays empty, so a positive budget is never reached. Algorithm 1’s cut counter also stays zero.\par
This is a complete execution argument for the printed rule, not a rounding issue or an assumption about data separability. It prevents the proposed starting procedure from producing its first cut. The conditional theorem assuming active cuts does not make a zero score negative.\par
Counterevidence is decisive: the release uses \(\texttt{np.sign(pred)}\ne y\), which treats zero as a mistake for labels \(\pm1\), and therefore can leave zero. Treating zero margins as eligible cuts restores this initial step; a cut through the centroid has a non-strict \(1-1/e\) bound. Neither full termination nor all other guarantees follows from this local change. Historical runs are not alleged to have stalled.\par
Problem 5 (Accepted cuts are counted instead of acquired labels).  Assessment: Major for the stated query-budget interpretation.\par
Algorithms 2/4 query a label before testing correctness, but add the point to \(D_{\rm AL}\) only if the cut is accepted. Their stopping condition counts \(|D_{\rm AL}|\). Algorithm 9’s baselines instead add every newly queried label. This matters directly to Section 7’s comparisons at a fixed number of queries.\par
The same accounting appears in `src/cpal/cpal.py`: the routine reads both `y[i\_mini]` and `y[i\_maxi]`, records tried indices, but increments `data\_used` only on mistakes. The sentiment application likewise plots and budgets `len(data\_used)`. Correctly classified queried labels have already been acquired, although they do not change the feasible set.\par
A complete local trace uses inputs \(1,2,3,4,5\), all labels \(+1\), both scalar patterns, and requested budget three. At the exact zero center, minimum and maximum scores tie at index zero. Both stale-center tests append that same index, giving `data\_used` length two but one distinct label. After these positive cuts, any strict barrier center has positive score on every remaining positive input. The next candidate pairs are indices \((1,4)\) and \((2,3)\); all four labels are inspected and all predictions are correct. Five distinct labels have now been acquired while the counter remains two, below the requested budget three. No numerical solver value is needed beyond the strictly positive cut slack.\par
The list also allows duplicate initial entries, so it is neither a distinct-label counter nor a reliable oracle-call counter. Fixing only the duplicate does not fix the uncharged correct labels. This demonstrates a core mismatch between the advertised label budget and the proposed accounting. It does not determine the actual number of labels behind the published curves or prove their ranking would reverse. A separate ledger of all queried distinct indices, with each label charged when requested, restores budget semantics; historical comparisons require those logs or fresh experiments. There is no certified numerical correction to the reported curves.\par
\section*{Additional released-code findings}
Problem 6 (Three-layer prediction reuses the inner-loop index).  Assessment: Major for the frozen release’s core implementation; historical impact unresolved.\par
In `src/cpal/utils.py`, `predict` iterates over outer index \(l\) and inner index \(j\), then multiplies both outer sums by `Di[i,j]` instead of `Di[i,l]`. The three-layer query and mistake tests call this routine, while the cut uses the intended outer pattern.\par
Take \(d=P_1=1\), \(P_2=2\), \(x=1\), \(Q=[1]\), and the generator-compatible sorted outer row \(D=[0,1]\). In the routine’s four parameter blocks choose \[c_1=(.1,.3),\ c_2=(.4,.1),\
 c_3=(.1,.2),\ c_4=(.2,.1).\] Their squared norm is \(.37<1\). All inner weights are positive; the outer preactivations are \((-.3,.2)\) and \((-.1,.1)\), strictly matching \(D\). The intended score is \(.2-.1=.1\). The released expression instead uses \(D_0=0\) for both outer terms and returns zero. This is an admissible sampled-pattern instance, not a violation of an activation constraint. Changing the selected index to \(l\) restores this displayed score and the corresponding general sum. The mathematical reformulation itself is not refuted; historical three-layer figures are not bound to this preliminary implementation.\par
This is Major for the frozen implementation’s core three-layer algorithm. In `src/cpal/cpal\_3layer.py`, line 55 uses this predictor to rank query candidates, and lines 78–83 use it to decide whether to cut and count a sample. The cut in lines 40–43 uses the intended outer pattern, so these central decisions use a different score from the constraint. The preliminary label and missing historical run binding limit attribution to the paper’s figures; they do not remove this demonstrated core implementation failure. A historical benchmark loss is not required to establish the defect in the inspected algorithm.\par
Problem 7 (Sentiment evaluation decodes the lifted vector in the wrong order).  Assessment: Major for the frozen release’s application evaluation; historical impact unresolved.\par
The sentiment application’s query and cut use interleaved pattern-major blocks through a Kronecker product. Its evaluation instead applies default row-major \(\texttt{c.reshape((2*d,m))}\). The required matrix layout has these blocks as columns.\par
For \(d=m=2\), take data rows \((-1,0)\) and \((1,0)\) with sorted pattern columns \[D=\begin{pmatrix}0&1\\1&0\end{pmatrix},\qquad
 c=(.3,0,.1,0,-.1,.8,-.2,0).\] Its squared norm is \(.79<1\). The intended positive and negative weight matrices are \[U_+=\begin{pmatrix}.3&-.1\\0&.8\end{pmatrix},
 \qquad
 U_-=\begin{pmatrix}.1&-.2\\0&0\end{pmatrix}.\] All signed preactivations on both data rows are strictly positive. At \((1,0)\), the correctly decoded ReLU score is \(.3-.1=.2\). The default reshape yields positive-block rows \((.3,0),(.1,0)\) and negative-block rows \((-.1,.8),(-.2,0)\), producing \(.3-.8=-.5\). Thus evaluation reverses the sign of a valid parameter vector. Column-major decoding restores the specified block correspondence.\par
The release also crops embeddings to 200 dimensions, whereas the paper reports 2560, and uses a different fixed data seed. These differences further prevent attributing Figure 4’s historical values to this exact code. The demonstrated evaluation defect remains; its historical score effect requires provenance and reruns.\par
This is Major for the frozen sentiment application’s evaluation path. Lines 168–172 compute and display train/test predictions with the incorrect layout, and lines 175–181 append that accuracy to the returned `score` sequence. The reversed classifier sign therefore reaches the application’s actual evaluation output. Query and cut calculations continue to use the separate Kronecker layout; they are not claimed to share this reshape defect. Historical uncertainty does not repair the demonstrated failure of this core evaluation path.\par
\section*{Overall dependence, surviving results, and empirical limits}
Problems 1–2 directly affect the advertised convergence theory. Problems 3–4 affect the specified training/active-learning algorithms, and Problem 5 affects the label-efficiency interpretation. They support Major without requiring a collapse of the convex reformulation or observed prediction scores. Problems 6–7 are also Major for the frozen release’s core three-layer decisions and sentiment evaluation, respectively. Their effect on historical paper figures remains unresolved because the exact release is not bound to those runs.\par
The reported quadratic-regression reductions in Table 5 are arithmetically consistent: \(1-.0111/.0568=80.4577\%\) for training and \(1-.0100/.0408=75.4902\%\) for test RMSE. These checks do not validate equal label costs. The spiral and regression studies use 100 points with an 80/20 split and five seeds for error bars. The IMDB experiment uses 50 training and 20 test examples, explicitly disclosed in Appendix H.1. It demonstrates a small selected experiment, not full-dataset validation.\par
The paper discloses sampled activation patterns, analytic-center optimization, and choosing whether to use a final regularized solve. Those choices must remain distinct from exact all-pattern and centroid theory. Baseline tuning, complete seed-level outputs, actual oracle histories, and the relationship between the preliminary code and the frozen figures remain unverified. No empirical rerun or external theorem reproduction was performed. No later correction was established from the available sources.\par
Item | Impact | Repair verification | Adoption / remaining work\par
Prediction convergence | Major | Diameter condition suffices | Unknown; algorithm premise\par
Inexact bounds | Major | Bounds refuted | New theorem required\par
Unit margin/radius | Major | Incompatibility proved | Valid radius or changed task\par
Zero initialization | Major | Local non-strict trigger | Release differs; full analysis\par
Label accounting | Major | Separate query ledger | Historical logs or reruns\par
Three-layer index | Major | Correct index verified | Release provenance needed\par
Sentiment reshape | Major | Block order verified | Historical scores unverified\par
Scientific adjudication is pending and `human\_reviews` is empty. Author adoption has not been requested. No claim of acceptance, human endorsement, or complete paper correctness is made.\par
9 E. Zhang, F. Zhang, and M. Pilanci. Active Learning of Deep Neural Networks via Gradient-Free Cutting Planes. ICML, PMLR 267 (2025). https://proceedings.mlr.press/v267/zhang25db.html. Staged 49-page proceedings PDF. Pilanci Lab. cpal, preliminary released repository. https://github.com/pilancilab/cpal. Commit `14c27c1a97756493738390775c1e4fe818d52935`; statically inspected source-only archive.\par
\section*{Supporting checks}
Reviewer-written `check.py` records the sector, ReLU volume, margin, query-count, decoder, and Table 5 calculations using existing Python and NumPy. The analytic G.2 construction covers every \(\Theta_4>0\); its numerical illustration does not assign the paper’s unspecified constant.\par
No author code was executed or additional download performed. External supplement discovery was unavailable; embedded appendices were available.\par
\end{document}
